\section{Тестирование и результаты}

В четвёртом разделе описывается проверка разработанного программного комплекса на реальных данных Яндекс.Недвижимость, оценка качества сбора и подготовки датасетов, а также пример прикладного анализа рынка через дашборд.

\subsection{Методика проверки}

Проверка проводилась без отдельного фреймворка автоматизированных UI-тестов: для учебного проекта достаточно сочетания контроля качества данных (QA-пайплайн), ручных сценариев по ролям и наблюдения за поведением системы на реальных выборках Москвы и Саратова.

\textbf{Уровни проверки:}
\begin{enumerate}
    \item \textbf{Данные парсера} --- стадии links, details, ETL; отчёт QA по контракту схемы; сверка числа строк после ingest на сервер.
    \item \textbf{Серверное API} --- авторизация, bootstrap датасета, фильтрация и экспорт CSV; запуск и мониторинг задач парсера из админ-панели.
    \item \textbf{Клиентский дашборд} --- загрузка серверного снимка и локального CSV; фильтры, KPI, графики, карта, таблица; разграничение возможностей по ролям.
    \item \textbf{Развёртывание} --- production-сборка, доступ по публичному URL (туннель или VPS), health-check API, работа карт с Referer production-хоста.
    \item \textbf{Интеграции} --- JavaScript API Яндекс.Карт на клиенте, API Геокодера и кэш SQLite на сервере (см.~рис.~\ref{fig:yandex-api}).
\end{enumerate}

\textbf{Учётные записи для сценариев} (seed при первом запуске сервера): \texttt{observer}, \texttt{analyst}, \texttt{admin}, \texttt{manager} с паролями, заданными в документации проекта. Каждая роль проходила типовой сценарий из табл.~\ref{tab:roles} (раздел~2).

\textbf{Критерии успешности:}
\begin{itemize}
    \item датасет города загружается в браузер без ошибок и отображает тысячи строк (Москва) или демонстрационную выборку (Саратов);
    \item фильтры меняют KPI и графики без заметной блокировки интерфейса;
    \item экспорт отфильтрованного CSV совпадает с таблицей на экране;
    \item администратор может запустить стадию парсера и увидеть прогресс; после ETL --- выполнить ingest;
    \item QA не фиксирует отсутствие обязательных колонок в processed CSV.
\end{itemize}

\subsection{Проверка пайплайна сбора и качества данных}

\subsubsection{Сценарий сбора для Саратова}

Для Саратова выполнен полный цикл: \texttt{links} $\rightarrow$ \texttt{details} $\rightarrow$ \texttt{etl} с ограничением \texttt{target\_listings=100}. По логам runtime-статуса: собрано порядка 1174 ссылок, обработано 95 карточек, после ETL --- 95 строк в \texttt{processed\_apartment\_data.csv}. Стадия завершилась без ошибок парсера; ingest добавил записи в серверный JSON-датасет с дедупликацией по ID объявления.

При повторном append часть строк была пропущена как дубликаты (в meta зафиксированы \texttt{skippedDuplicates}), что подтверждает работу слияния при загрузке.

\subsubsection{Сценарий для Москвы}

Московский датасет формировался несколькими партиями append из processed CSV. Итоговый серверный снимок содержит \textbf{12\,549} объявлений (обновление от 25.05.2026). В meta сохранены две загрузки: первая партия --- 5341 строка, вторая --- 7208 добавленных при 3217 пропущенных дубликатах и 16 отфильтрованных как слишком удалённых от центра города.

\subsubsection{Отчёт QA}

Модуль \texttt{pipeline/qa\_report.py} проверяет наличие \texttt{required\_core} колонок (цена, комнаты, площадь, координаты), долю пропусков (предупреждение при $>$15\%), дубликаты ссылок. Для сданных датасетов критических ошибок схемы не выявлено; предупреждения возможны по отдельным полям с высокой долей пропусков (например, текстовый адрес, если координаты есть, но reverse geocode ещё не применён).

\subsection{Результаты на данных Москвы и Саратова}

Сводные показатели серверных снимков приведены в табл.~\ref{tab:results-cities}.

\begin{table}[htbp]
\centering
\small
\begin{tabular}{|L{2.4cm}|L{2.6cm}|L{2.8cm}|L{2.8cm}|L{2.8cm}|}
\hline
\textbf{Город} & \textbf{Число объявлений} & \textbf{Медиана цены, млн руб.} & \textbf{Медиана площади, м$^2$} & \textbf{Медиана расст. до центра, км} \\
\hline
Москва & 12\,549 & 24,0 & 55,2 & 11,0 \\
\hline
Саратов & 75 & 4,1 & 39,8 & 7,4 \\
\hline
\end{tabular}
\caption{Сводные характеристики серверных датасетов (по состоянию на 25.05.2026)}
\label{tab:results-cities}
\end{table}

\textbf{Москва} --- основная выборка для нагрузочной проверки клиента: bootstrap JSON $\approx$47~МБ, фильтрация и перерисовка графиков на десятках тысяч строк остаются приемлемыми за счёт \texttt{useMemo} и \texttt{useDeferredValue} (рис.~\ref{fig:dashboard-main}). Карта с лимитом 2000 точек по умолчанию показывает кластеры по агломерации (рис.~\ref{fig:dashboard-map}); увеличение лимита до нескольких тысяч возможно, но увеличивает время отрисовки.

\textbf{Саратов} --- демонстрационный регион с меньшим объёмом данных; удобен для быстрой проверки полного цикла парсера и ingest без длительного ожидания details.

\subsubsection{Проверка развёртывания}

Помимо локального \texttt{npm run dev:all}, система проверялась в двух сетевых конфигурациях. \textbf{Демонстрационный туннель} (\texttt{npm run dev:tunnel}): production-бандл и API пробрасываются наружу через localtunnel; подтверждены вход по JWT, загрузка bootstrap Москвы и отображение карты при корректном Referer. \textbf{VPS} (Ubuntu на Beget): выполнена установка скриптом \texttt{deploy/setup.sh} --- nginx на порту~80, unit \texttt{real-estate-dashboard}, статика из \texttt{dist/}, проксирование \texttt{/api} на \texttt{127.0.0.1:3001}. После деплоя успешно проходят \texttt{GET /api/health}, открытие главной страницы и сценарии analyst/admin из табл.~\ref{tab:roles-test}. Схема production-размещения --- рис.~\ref{fig:deploy-vps} (раздел~3).

\subsection{Проверка сценариев по ролям}

Результаты ручной проверки согласованы с матрицей доступа (табл.~\ref{tab:roles-test}).

\begin{table}[htbp]
\centering
\small
\begin{tabular}{|L{2.2cm}|L{4.2cm}|L{7.0cm}|}
\hline
\textbf{Роль} & \textbf{Сценарий} & \textbf{Результат} \\
\hline
observer & Вход, Москва, фильтры, KPI, графики, карта & Доступен только серверный датасет; таблица и экспорт недоступны; упрощённый список объявлений отображается \\
\hline
analyst & Сервер + свой CSV, конструктор графиков, экспорт & Переключение источника работает; CSV до 50\,000 строк загружается чанками; экспорт CSV/XLSX совпадает с фильтром \\
\hline
admin & Upload, ingest, запуск парсера, прогресс & Ingest processed CSV обновляет meta; parser job показывает стадию и процент; resume после капчи принимается через API \\
\hline
manager & Вкладки «Рынки» и «Пользователи» & Сводка по городам совпадает с meta; создание пользователя и смена роли сохраняются в SQLite \\
\hline
\end{tabular}
\caption{Результаты проверки типовых сценариев по ролям}
\label{tab:roles-test}
\end{table}

Проверка роли \texttt{observer} подтвердила доступ только к серверному датасету и отсутствие таблицы с экспортом (рис.~\ref{fig:observer-dashboard}). Для \texttt{analyst} работают переключение источника, конструктор графиков и выгрузка --- на рис.~\ref{fig:dashboard-table}. Администратор успешно выполнял ingest (рис.~\ref{fig:admin-server}) и отслеживал задачи парсера (рис.~\ref{fig:admin-parser}). Руководитель видит сводку по рынкам (рис.~\ref{fig:manager-markets}) и управляет учётными записями (рис.~\ref{fig:manager-users}).

\textbf{Авторизация:} неверный пароль возвращает HTTP 401; без токена защищённые маршруты недоступны. JWT действует 12~часов, после истечения требуется повторный вход.

\textbf{Согласованность фильтров:} при одинаковых параметрах query клиентская таблица, KPI и \texttt{GET /api/listings} дают одинаковое число строк (проверено на выборке «двушки, цена до 15~млн, расстояние до 10~км» для Москвы).

\subsection{Пример прикладного анализа}

Рассмотрим сегмент вторичного рынка двухкомнатных квартир в Москве в радиусе 10~км от центра с ценой до 15~млн руб. (типичный сценарий из раздела~2). Состояние дашборда после применения фильтров показано на рис.~\ref{fig:analysis-filters}.

\screenshot{analysis-filters}{analysis-filters}{Прикладной срез: двухкомнатные квартиры, цена до 15~млн руб., расстояние до центра до 10~км (Москва)}

После применения фильтров в выборке остаётся порядка нескольких сотен объявлений (точное число зависит от актуальности снимка). Медиана цены оказывается ниже общегородской медианы (24~млн), медиана площади --- около 45--55~м$^2$. На гистограмме цены (рис.~\ref{fig:analysis-charts}) виден левый сдвиг распределения; scatter «цена vs площадь» показывает положительную связь с разбросом на уровне старых панельных серий.

\screenshot{analysis-charts}{analysis-charts}{Гистограмма цены и scatter «цена~--- площадь» для отфильтрованного сегмента}

Блок \textbf{«Сравнение сегментов»} позволяет рядом поставить, например, «новостройки» (фильтр по году постройки $\geq$ 2015) и оставшуюся вторичку (рис.~\ref{fig:analysis-segments}): медиана цены за м² у новостроек, как правило, выше, что согласуется с общими представлениями о премиальности первичного рынка.

\screenshot{analysis-segments}{analysis-segments}{Сравнение новостроек (год $\geq$ 2015) и вторичного рынка в том же географическом срезе}

На \textbf{карте} объекты сегмента группируются в кольце вокруг центра (рис.~\ref{fig:dashboard-map}); балуны содержат цену и ссылку на карточку. Для строк без адреса выполнялось обратное геокодирование: часть записей получила текстовый адрес через кэш SQLite после нескольких запросов к API.

\textbf{Экспорт:} отфильтрованная таблица выгружена в CSV и сверена с экраном --- расхождений по числу строк нет.

\subsection{Производительность и ограничения}

При работе с московским снимком (12\,549 строк) первичная загрузка bootstrap занимает несколько секунд на локальном сервере; последующие изменения фильтров не вызывают повторных запросов к API. Импорт CSV аналогичного размера на клиенте выполняется с паузами между чанками --- интерфейс остаётся отзывчивым.

Выявленные \textbf{ограничения}:
\begin{itemize}
    \item парсер чувствителен к капче Яндекса --- требуется ручное вмешательство или headless$\rightarrow$окно;
    \item не все объявления содержат полный адрес; геокодирование ограничено квотами API;
    \item SQLite и JSON-файлы рассчитаны на локальное/учебное развёртывание, не на высокую конкурентную нагрузку;
    \item саратовская выборка мала для устойчивых выводов о структуре рынка --- годится для отладки, не для статистической репрезентативности.
\end{itemize}

\subsection{Выводы по четвёртому разделу}

Проведённая проверка показала, что программный комплекс выполняет заявленные функции:
\begin{enumerate}
    \item пайплайн Selenium собирает и нормализует данные; QA и ingest обеспечивают контролируемое попадание строк на сервер;
    \item серверные датасеты Москвы (12\,549 объявлений) и Саратова (75 объявлений) доступны через API и дашборд;
    \item ролевая модель ограничивает операции записи и экспорт в соответствии с проектом;
    \item интерактивный анализ (фильтры, графики, карта, сравнение сегментов) применим к реальным данным и поддерживает прикладные срезы рынка;
    \item развёртывание на VPS и демонстрация через HTTPS-туннель подтверждают пригодность комплекса для показа вне локальной машины разработчика.
\end{enumerate}

Система пригодна для учебной и исследовательской работы с объявлениями Яндекс.Недвижимость; для промышленной эксплуатации потребовались бы масштабирование хранилища, автоматизация UI-тестов и более устойчивый сбор данных в условиях anti-bot защиты.
